% Transcribed human main text; notation and source in tex/010.
\begin{lemma}[15.7]
Let $c$ be an Arakelov divisor of a global field $K$. The set $L(c)$ is finite.
\end{lemma}
\begin{proof}
We assume $K$ is a number field; see Problem Set 7 for the function field case. The fractional ideal $I_c$ is a lattice in $K_{\mathbb R}$ (under the canonical embedding $K\hookrightarrow K_{\mathbb R}$), and is thus a closed discrete subset of $K_{\mathbb R}$ (recall from Remark 14.5 that lattices are closed). In $K_{\mathbb R}$ we may view $L(c)=I_c\cap R_c$ as the intersection of a discrete closed set with a compact set, which is a compact discrete set and therefore finite.
\end{proof}
\begin{corollary}[15.8]
Let $K$ be a global field, and let $\mu_K$ denote the torsion subgroup of $K^\times$ (equivalently, the roots of unity in $K$). The group $\mu_K$ is finite and equal to the kernel of the map $K^\times\to\operatorname{Div}K$ defined by $x\mapsto(\|x\|_v)$; it is also the torsion subgroup of $\mathcal O_K^\times$.
\end{corollary}
\begin{proof}
Each $\zeta\in\mu_K$ satisfies $\zeta^n=1$ for some positive integer $n$. For every place $v\in M_K$ we have $\|\zeta^n\|_v=\|\zeta\|_v^n=1$, and therefore $\|\zeta\|_v=1$. It follows that $\mu_K\subseteq\ker(K^\times\to\operatorname{Div}K)$. Let $c$ be the Arakelov divisor with $c_v=1$ for all $v\in M_K$. Then $\ker(K^\times\to\operatorname{Div}K)\subseteq L(c)$ is a finite subgroup of $K^\times$ and is therefore contained in the torsion subgroup $\mu_K$. Every element of $\mu_K$ is an algebraic integer (in fact a root of $x^n-1$), so $\mu_K\subseteq\mathcal O_K^\times$.
\end{proof}

\begin{proposition}[15.9]
Let $K$ be a number field with $s$ complex places, define
\[
B_K:=\left(\frac2\pi\right)^s\sqrt{|D_K|}.
\]
If $c$ is an Arakelov divisor of size $\|c\|>B_K$ then $L(c)$ contains an element of $K^\times$.
\end{proposition}
\begin{proof}
Our strategy is to apply Minkowski's lattice point theorem (see Theorem 14.13) to the convex symmetric set $R_c$ and the lattice $I_c\subseteq K\subseteq K_{\mathbb R}$; we just need to show that if $\|c\|>B_K$ then the ratio of the Haar measure of $R_c$ to the covolume of $I_c$ exceeds $2^n$, where $n=r+2s$ is the degree of $K$ (which is the real dimension of $K_{\mathbb R}$). As defined in \S14.2, we normalize the Haar measure $\mu$ on the locally compact group $K_{\mathbb R}\simeq\mathbb R^r\times\mathbb C^s\simeq\mathbb R^n$ so that $\mu(S)=2^s\mu_{\mathbb R^n}(S)$ for measurable $S\subseteq K_{\mathbb R}$. For each real place $v$, the constraint $\|x\|_v=|x|_{\mathbb R}\leq c_v$ contributes a factor of $2c_v$ to $\mu(R_c)$, and for each complex place $v$ the constraint $\|x\|_v=|x|_{\mathbb C}^2\leq c_v$ contributes a factor of $\pi c_v$ (the area of a circle of radius $\sqrt{c_v}$). We may then compute
\begin{align*}
\frac{\mu(R_c)}{\covol(I_c)}
&=\frac{2^s\mu_{\mathbb R^n}(R_c)}{\covol(I_c)}
=\frac{2^s\bigl(\prod_{v\text{ real}}2c_v\bigr)\bigl(\prod_{v\text{ complex}}\pi c_v\bigr)}{\covol(I_c)}\\
&=\frac{2^r(2\pi)^s\prod_{v\mid\infty}c_v}{\sqrt{|D_K|}N(I_c)}
=\frac{2^r(2\pi)^s}{\sqrt{|D_K|}}\|c\|
=\frac{\|c\|}{B_K}2^n>2^n
\end{align*}
where we have used Corollary 14.17 and (1) in the second line. Theorem 14.13 implies that $L(c)=R_c\cap I_c$ contains a nonzero element (which lies in $K^\times\subseteq K_{\mathbb R}$, since $I_c\subseteq K\subseteq K_{\mathbb R}$).
\end{proof}

\begin{proposition}[15.11]
Let $K$ be a number field with $r$ real and $s$ complex places, and let $\Lambda_K$ be the image of the unit group $\mathcal O_K^\times$ in $\mathbb R^{r+s}_0$ under the $\operatorname{Log}$ map. The following hold:
\begin{enumerate}
\item We have a split exact sequence of finitely generated abelian groups
\[
1\longrightarrow\mu_K\longrightarrow\mathcal O_K^\times
\xrightarrow{\operatorname{Log}}\Lambda_K\longrightarrow0;
\]
\item $\Lambda_K$ is a lattice in the trace zero hyperplane $\mathbb R^{r+s}_0$.
\end{enumerate}
\end{proposition}
Here $\mu_K$ is not a Haar measure, it denotes the group of roots of unity in $K$, all of which are clearly torsion elements of $\mathcal O_K^\times$, and any torsion element of $\mathcal O_K^\times$ is clearly a root of unity.
\begin{proof}
(1) We first show exactness. Let $Z$ be the kernel of $\mathcal O_K^\times\xrightarrow{\operatorname{Log}}\Lambda_K$. Clearly $\mu_K\subseteq Z$, since $\Lambda_K\subseteq\mathbb R^{r+s}_0$ is torsion free. Let $c$ be the Arakelov divisor with $I_c=\mathcal O_K$ and $c_v=2$ for $v\mid\infty$, so that
\[
L(c)=\{x\in\mathcal O_K:\|x\|_v\leq2\text{ for all }v\mid\infty\}.
\]
For $x\in\mathcal O_K^\times$ we have
\[
x\in L(c)\ \Longleftrightarrow\ \operatorname{Log}(x)\in\operatorname{Log}R_c
=\{z\in\mathbb R^{r+s}:z_i\leq\log2\}.
\]
The set on the RHS includes the zero vector, thus $Z\subseteq L(c)$, which by Lemma 15.7 is a finite set. As a finite subgroup of $\mathcal O_K^\times$, we must have $Z\subseteq\mu_K$, so $Z=\mu_K$ and the sequence is exact (the map from $\mathcal O_K^\times$ to $\Lambda_K$ is surjective by the definition of $\Lambda_K$).

We now show the sequence splits. Note that
$\Lambda_K\cap\operatorname{Log}(R_c)=\operatorname{Log}(\mathcal O_K^\times\cap L(c))$ is finite, since $L(c)$ is finite. It follows that $0$ is an isolated point of $\Lambda_K$ in $\mathbb R^{r+s}$, and in $\mathbb R^{r+s}_0$, so $\Lambda_K$ is a discrete subgroup of $\mathbb R^{r+s}_0\simeq\mathbb R^{r+s-1}$ and must be finitely generated, by Lemma 14.3. It follows that $\mathcal O_K^\times$ is finitely generated, since it lies in a short exact sequence whose left and right terms are finitely generated (recall that $\mu_K$ is finite, by Corollary 15.8). By the structure theorem for finitely generated abelian groups, the sequence must split, since $\mu_K$ is the torsion subgroup of $\mathcal O_K^\times$.

(2) Having proved (1) it remains only to show that $\Lambda_K$ spans $\mathbb R^{r+s}_0$. Let $V$ be the subspace of $\mathbb R^{r+s}_0$ spanned by $\Lambda_K$ and suppose for the sake of contradiction that $\dim V<\dim\mathbb R^{r+s}_0$. The orthogonal subspace $V^\perp$ then contains a unit vector $u$, and for every $\lambda\in\mathbb R_{>0}$ the open ball $B_{<\lambda}(\lambda u)$ does not intersect $\Lambda_K$. Thus $\mathbb R^{r+s}_0$ contains points arbitrarily far away from every point in $\Lambda_K$ (with respect to any norm on $\mathbb R^{r+s}_0\subseteq\mathbb R^{r+s}$). To obtain a contradiction it is enough to show that there is a constant $M\in\mathbb R_{>0}$ such that for every $h\in\mathbb R^{r+s}_0$ there is an $\ell\in\Lambda_K$ for which $\|h-\ell\|:=\max_i|h_i-\ell_i|<M$ (here we are using $\|\ \|$ to denote the sup norm on the $\mathbb R$-vector space $\mathbb R^{r+s}$).

Let us fix a real number $B>B_K$, where $B_K$ is as in Proposition 15.9, so that for every $c\in\operatorname{Div}K$ with $\|c\|\geq B$ the set $L(c)$ contains a nonzero element, and fix a vector $b\in\mathbb R^{r+s}$ with nonnegative components $b_i$ such that $T(b)=\sum_i b_i=\log B$. Let $(\alpha_1),\ldots,(\alpha_m)$ be the list of all nonzero principal ideals with $N(\alpha_j)\leq B$ (by Lemma 14.21 this is a finite list). Let $M$ be twice the maximum of $(r+s)B$ and $\max_j\|\operatorname{Log}(\alpha_j)\|$.

Now let $h\in\mathbb R^{r+s}_0$, and define $c\in\operatorname{Div}K$ by $I_c:=\mathcal O_K$ and $c_v:=\exp(h_i+b_i)$ for $v\mid\infty$, where $i$ is the coordinate in $\mathbb R^{r+s}$ corresponding to $v$ under the $\operatorname{Log}$ map. We have
\begin{align*}
\|c\|=\prod_v c_v
&=\exp\left(\sum_i(h_i+b_i)\right)
=\exp T(h+b)\\
&=\exp(T(h)+T(b))=\exp T(b)=B>B_K,
\end{align*}
thus $L(c)$ contains a nonzero $\gamma\in I_c\cap K=\mathcal O_K$, and $g=\operatorname{Log}(\gamma)$ satisfies $g_i\leq\log c_v=h_i+b_i$. We also have $T(g)=T(\operatorname{Log}\gamma)=\log N(\gamma)\geq0$, since $N(\gamma)\geq1$ for all nonzero $\gamma\in\mathcal O_K$. The vector $w:=g-h\in\mathbb R^{r+s}$ satisfies $\sum_i w_i=T(v)=T(g)-T(h)=T(g)\geq0$ and $w_i\leq b_i\leq B$ which together imply $|w_i|\leq(r+s)B$, so $\|g-h\|=\|w\|\leq M/2$. We also have
\[
\log N(\gamma)=T(\operatorname{Log}(\gamma))\leq T(h+b)=T(b)=\log B,
\]
so $N(\gamma)\leq B$ and $(\gamma)=(\alpha_j)$ for one of the $\alpha_j$ fixed above. Thus $\gamma/\alpha_j\in\mathcal O_K^\times$ is a unit, and
\[
\ell:=\operatorname{Log}(\gamma/\alpha_j)=\operatorname{Log}(\gamma)-\operatorname{Log}(\alpha_j)\in\Lambda_K
\]
satisfies $\|g-\ell\|=\|\operatorname{Log}(\alpha_j)\|\leq M/2$. We then have
\[
\|h-\ell\|\leq\|h-g\|+\|g-\ell\|\leq M
\]
as desired (by the triangle inequality for the sup-norm).
\end{proof}

Dirichlet's unit theorem follows immediately from Proposition 15.11.
\begin{theorem}[15.12 (Dirichlet's Unit Theorem)]
Let $K$ be a number field with $r$ real and $s$ complex places. Then $\mathcal O_K^\times\simeq\mu_K\times\mathbb Z^{r+s-1}$ is a finitely generated abelian group.
\end{theorem}
\begin{proof}
The image of the torsion-free part of the unit group $\mathcal O_K^\times$ under the $\operatorname{Log}$ map is the lattice $\Lambda_K$ in the trace-zero hyperplane $\mathbb R^{r+s}_0$, which has dimension $r+s-1$.
\end{proof}
